\subsection{Integrals of operators}

Let G and H be two Hausdorff locally convex spaces over R, and suppose now that F is the space \(\mathcal{L}(\mathrm{G};\mathrm{H})\) of continuous linear mappings of G into H, equipped with the topology of pointwise convergence. Every continuous linear form on F may be extended to a continuous linear form on the product space \(H^{G}\) (TVS, II, §4, No.~1, Prop. 2), hence may be written \(u \mapsto \sum_{i=1}^{n} \langle u(\mathbf{a}_{i}), \mathbf{b}_{i}^{\prime} \rangle\) , where the \(a_{i}\) (resp. the \(b_{i}^{\prime}\) ) are elements of G (resp. of the dual \(H^{\prime}\) of H). To say that a mapping U of T into F is scalarly essentially \(\mu\) -integrable means that, for every \(a \in G\) and every \(b \in H^{\prime}\) , the numerical function \(t \mapsto \langle U(t) \cdot a, b^{\prime} \rangle\) is essentially \(\mu\) -integrable.

PROPOSITION 9. --- Let U be a scalarly essentially \(\mu\) -integrable mapping of T into \(\mathbf{F} = \mathcal{L}_{s}(\mathbf{G}; \mathbf{H})\) . In order that \(\int U d\mu \in F\) , it is necessary and sufficient that the following two conditions be satisfied:

\begin{enumerate}
\def\labelenumi{\alph{enumi})}
\item
  For every \(\mathbf{x} \in \mathbf{G}\) , \(\int (U(t) \cdot \mathbf{x}) d\mu(t) \in \mathrm{H}\) .
\item
  For every equicontinuous subset \(B'\) of \(H'\) , the set of linear forms \(u_{y'} : x \mapsto \int \langle U(t) \cdot x, y' \rangle d\mu(t)\) , where \(y'\) runs over \(B'\) , is equicontinuous.
\end{enumerate}

The conditions a) and b) are necessary. For, since for every \(\mathbf{x} \in G\) , the mapping \(\widetilde{\mathbf{x}}: V \mapsto V \cdot \mathbf{x}\) of \(\mathcal{L}_s(\mathrm{G};\mathrm{H})\) into H is linear, one sees (No.~1,

Prop. 1) that \(\widetilde{\mathbf{x}}\circ U:t\mapsto U(t)\cdot \mathbf{x}\) is scalarly essentially \(\mu\) -integrable and that

\[
S \cdot \mathbf {x} = \int (U (t) \cdot \mathbf {x}) d \mu (t),\tag{1}
\]

where \(S = \int U d\mu \in \mathcal{L}_s(\mathrm{G};\mathrm{H})\) . This proves a). Moreover, (1) may also be written

\[
\langle S \cdot \mathbf {x}, \mathbf {y} ^ {\prime} \rangle = \int \left\langle U (t) \cdot \mathbf {x}, \mathbf {y} ^ {\prime} \right\rangle d \mu (t) = \left\langle \mathbf {x}, u _ {\mathbf {y} ^ {\prime}} \right\rangle ,\tag{2}
\]

in other words \(^{t}S \cdot y' = u_{y'}\) . Since S is continuous, \(^{t}S\) transforms every equicontinuous subset of \(H'\) into an equicontinuous subset of \(G'\) , whence b).

Conversely, suppose a) and b) verified. By virtue of a), the formula (1) defines a linear mapping S of G into H, and, for every \(y' \in H'\) , this mapping satisfies (2) (No.~1, Prop. 1); but then, condition (b) says that S is continuous (TVS, II, §6, No.~4, Props. 5 and 6, and III, §3, No.~5, Prop. 7), therefore \(S \in \mathcal{L}_{s}(G;H)\) . Finally, formula (2) proves that \(S = \int U d\mu\) .

COROLLARY. --- The condition b) of Prop. 9 is satisfied in each of the following two cases:

\(1^{\circ}\) The measure \(\mu\) is bounded, and if S is its support, then \(U(S)\) is an equicontinuous subset of \(\mathcal{L}(G;H)\) .

\(2^{\circ}\) The condition a) of Prop. 9 is satisfied, the space G is barreled, and, for every compact subset K of T, \(U(\mathrm{K})\) is a bounded subset of \(\mathcal{L}_s(\mathrm{G};\mathrm{H})\) .

Let us first place ourselves in case \(1^{\circ}\) . We can restrict ourselves to the case that S = T (Ch. V, §7, No.~1, Th. 1). Then, for every equicontinuous subset \(B'\) of \(H'\) , there exists an equicontinuous, convex, balanced and weakly closed subset \(A' \subset G'\) such that \({}^{t}U(t) \cdot y' \in A'\) for all \(y' \in B'\) and all \(t \in T\) (TVS, II, §6, No.~4, Prop. 6). Since U is scalarly \(\mu\) -integrable, the mapping \(t \mapsto {}^{t}U(t) \cdot y'\) of T into the dual \(G'\) of G equipped with \(\sigma(G', G)\) , is scalarly \(\mu\) -integrable, and one may write

\[
u _ {\mathbf {y} ^ {\prime}} = \int \left(^ {t} U (t) \cdot \mathbf {y} ^ {\prime}\right) d \mu (t).
\]

Since \(A'\) is convex and compact for \(\sigma(G', G)\) , the Cor. of Prop. 5 of No.~2 shows that \(u_{\mathbf{y}'} \in \mu(T)A'\) for every \(\mathbf{y}' \in B'\) , which proves our assertion.

Let us now place ourselves in case \(2^{\circ}\) . For every \(y' \in H'\) and every compact subset K of T, set

\[
u _ {\mathrm{K}, \mathbf {y} ^ {\prime}} = \int \varphi_ {\mathrm{K}} (t) \left(^ {t} U (t) \cdot \mathbf {y} ^ {\prime}\right) d \mu (t),
\]

an element of the algebraic dual \(\mathbf{G}^*\) of \(\mathbf{G}\) . Since \(\mathbf{G}\) is barreled, every bounded subset of \(\mathcal{L}_s(\mathbf{G};\mathbf{H})\) is equicontinuous (TVS, III, §4, No.~2, Th. 1); the first part of the argument, applied to the function \(\varphi_{\mathbf{K}}U\) and the bounded measure \(\varphi_{\mathbf{K}}\cdot \mu\) , shows that \(u_{\mathbf{K},\mathbf{y}'}\in \mathbf{G}'\) . Moreover, for the topology \(\sigma (\mathbf{G}^{*},\mathbf{G})\) , one has \(u_{\mathbf{y}'} = \lim_{\mathbf{K}} = u_{\mathbf{K},\mathbf{y}'}\) , the limit being taken with respect to the increasing directed set of compact subsets of \(\mathbf{T}\) (Ch. V, §1, No.~3, Prop. 10). To verify condition b) of Prop. 9, it suffices, by Prop. 9, to prove that the linear mapping \(S\) of \(\mathbf{G}\) into \(\mathbf{H}\) defined by (1) is continuous; moreover, \(\mathbf{G}\) is barreled, therefore it will suffice to prove that \(S\) is continuous when \(\mathbf{G}\) and \(\mathbf{H}\) are equipped with their weakened topologies (TVS, IV, §1, No.~3, Prop. 7); finally, by virtue of (2), one is reduced to showing that \(u_{\mathbf{y}'}\in \mathbf{G}'\) for every \(\mathbf{y}'\in \mathbf{H}'\) . Since \(u_{\mathbf{y}'}\) is in the closure, for \(\sigma (\mathbf{G}^{*},\mathbf{G})\) , of the set \(\mathbf{M}'\) of the \(u_{\mathbf{K},\mathbf{y}'}\) , where \(\mathbf{K}\) runs over the set of compact subsets of \(\mathbf{T}\) , it suffices to prove that \(\mathbf{M}'\) is equicontinuous (TVS, III, §3, No.~4, Prop. 4); and since \(\mathbf{G}\) is barreled, it comes to the same to say that for every \(\mathbf{x}\in \mathbf{G}\) , the set of \(\langle \mathbf{x},u_{\mathbf{K},\mathbf{y}'}\rangle\) is bounded (TVS, III, §4, No.~2, Th. 1). But this follows at once from the relations

\[
| \langle \mathbf {x}, u _ {\mathrm{K}, \mathbf {y} ^ {\prime}} \rangle | = \left| \int \varphi_ {\mathrm{K}} (t) \langle U (t) \cdot \mathbf {x}, \mathbf {y} ^ {\prime} \rangle d \mu (t) \right| \leqslant \int | \langle U (t) \cdot \mathbf {x}, \mathbf {y} ^ {\prime} \rangle | d \mu (t).
\]

PROPOSITION 10. --- Let U be a mapping of T into F = \(\mathcal{L}_s(\mathrm{G};\mathrm{H})\) . In each of the following three cases, U is scalarly essentially \(\mu\) -integrable and \(\int U d\mu \in \mathcal{L}_s(\mathrm{G};\mathrm{H})\) :

\begin{enumerate}
\def\labelenumi{\alph{enumi})}
\item
  H is quasi-complete, \(\mu\) is bounded and, if S is its support, U is \(\mu\) -measurable and \(U(S)\) is equicontinuous.
\item
  H is semi-reflexive, \(\mu\) is bounded and, if S is its support, U is scalarly \(\mu\) -measurable and \(U(S)\) is equicontinuous.
\item
  H is semi-reflexive, G is barreled, U is scalarly essentially \(\mu\) -integrable, and, for every compact subset K of T, \(U(\mathrm{K})\) is bounded.
\end{enumerate}

The fact that U is scalarly essentially integrable is obvious in all three cases; by virtue of Prop. 9 and its corollary, it suffices in each case to verify condition a) of Prop. 9. Now, this condition follows from Prop. 8 of No.~2 in the first case, and from the Cor. of Prop. 7 of No.~2 in the other two cases.

